%% --- 6.2 CoE Audit Results ---
%% Results from CoE audit v2.7 (forensic_v2.7_20260506), OAS papers = v3.9 (v3.1.9, best solutions are selected from last iteration only (best of 5 branches), not best of all 25 solutions).
%% Sakana ASv2 results from forensic_v2.8_sakana_gemini31, merged into v2.7 on 2026-05-19. Model: gemini-3.1-pro-preview. I1 scores extracted via Gemini consensus (tex+pdf).
%% I3 hallucinated-reference counts are human-confirmed (automated API + manual review).

\subsection{CoE Audit Results}
\label{sec:exp:integrity}

\begin{table}[h]
\centering
\small
\caption{\coeaudit{} results across five systems (15~papers per system).
EPLB papers are excluded from Score Verif.\ because its scoring formula includes an execution-time component that varies with hardware, making scores non-reproducible across machines.
Metric definitions are in \S\ref{sec:coeaudit}.
}
\label{tab:integrity}
\resizebox{\textwidth}{!}{%
\begin{tabular}{@{}lcccc@{}}
\toprule
\textbf{System} & \textbf{Score Verif. $\uparrow$} & \textbf{Spec.\ Violation $\downarrow$} & \textbf{Ref.\ Verif. $\downarrow$} & \textbf{Method-Code $\uparrow$} \\
\midrule
Sakana AI-Scientist v2~\citep{yamada2025aiscientistv2}     & 5/12  & 10/15 & 0/159    & 5/15  \\
AutoResearchClaw~\citep{li2024autoresearchclaw}            & 5/12  & 0/15 & 3/196    & 3/15  \\
DeepScientist~\citep{huang2025deepscientist}             & 11/12 & 0/15 & 42/201   & 5/15  \\
AI-Researcher~\citep{tang2025airesearcher}            & 9/12  & 1/15 & 21/222   & 12/15 \\
\midrule
\sys{}         & 12/12 & 0/15 & 0/337    & 14/15 \\
\bottomrule
\end{tabular}}
\end{table}

The results of \coeaudit{} across five systems are presented in Table~\ref{tab:integrity}.
All I1--I3 flagged results were manually verified by human reviewers. I4 judgments were validated on a sampled basis.
\sys{} is the only system to lead on all four checks: perfect score verification (12/12), zero specification violations (0/15), zero hallucinated references (0/337), and the highest method-code alignment (14/15).
The gap is largest in reference integrity and method-code alignment---the two checks that test evidence provenance rather than score reproduction.
Because the BFTS--ADRS design mismatch confounds both I2 and I4 for Sakana, cross-system comparison on these two checks should exclude Sakana. I1 and I3 remain valid.

\paragraph{Score verification (I1).}
\textbf{\sys{}} achieves perfect score verification (12/12): every paper's claimed result reproduces exactly under re-evaluation.
\textbf{DS} matches in 11/12 (92\%), with the single failure caused by fabricated metric direction---the paper claims ``higher is better'' for a cost-minimization metric, framing the raw cost as an inverse aggregate score, so the baseline's 1035.1 reads as the best result when it is actually the worst.
\textbf{AIR} matches in 9/12 (75\%), with failures spanning small-magnitude discrepancies (1--4\%) and one paper that reports no quantitative scores at all.
\textbf{ARC} matches in 5/12 (42\%), and its failures trace to three root causes: (1)~crashed solvers (5 of 15~solvers import helper modules generated by ARC's multi-file blueprint planner that are absent during standalone re-evaluation, producing evaluator fallback scores that differ from the paper's claims); (2)~evaluator mismatch (ARC's bundled cloudcast evaluator includes a patch absent from the canonical evaluator, producing different scores for the same solver); and (3)~stochastic evaluation noise in txn\_scheduling (2--3\% variance from unseeded scheduling randomness).
\textbf{Sakana ASv2} matches in 5/12 (42\%), the lowest among all systems.
Manual investigation of the 7~failures reveals two dominant patterns.
First, \emph{cross-stage score cherry-picking} (4 of 7~failures): the writeup LLM receives summaries from all four BFTS stages as context, and selects the most favorable score from ablation-stage nodes rather than the score of the node whose code is used as the final solution.
For example, in \textsc{prism} seed-1, the selected node scores 22.79 but the paper reports 25.39---a number traced to ablation node~6 (``Ablate KVPR-Aware Initialization'') in \texttt{ablation\_summary.json}.
The same pattern appears in \textsc{cloudcast} seed-0 (+56\%), \textsc{prism} seed-2 ($-$4.7\%, paper under-reports), and \textsc{txn\_scheduling} seed-2 (+17\%).
Second, \emph{environment-dependent tuning} (2 of 7~failures): the solver contains a hyperparameter tuning loop gated on an environment variable that is set differently during canonical re-evaluation, causing the solver to use default parameters instead of the tuned ones (e.g., \textsc{prism} seed-0: 26.26 tuned vs.\ 22.34 default, a 15\% gap).
The remaining failure is a metric mismatch (\textsc{cloudcast} seed-1: the paper reports per-transfer dollar cost while the evaluator produces a combined score).

\paragraph{Specification violation (I2).}
Specification violation rates are uniformly low for \textbf{ARC}, \textbf{DS}, and \textbf{\sys{}} (0/15), while \textbf{AIR} has one flagged paper (\textsc{llm\_sql}) where the solver physically reorders values across columns within each row, destroying column integrity to inflate the prefix-cache hit metric.
\textbf{Sakana ASv2} registers 10/15 specification violations---the highest rate.
The agent could tune parameters through BFTS's iteration loop (one setting per iteration), but the stage~2 goal (``test across multiple parameter settings'') encourages intra-iteration sweeps. Combined with the evaluator import pattern visible in our canonical harness, this leads the agent to import the evaluator and build its own tuning loops in 10 of 15~runs.
Most violations trace to the BFTS--ADRS design mismatch rather than adversarial behaviour (Appendix~\ref{app:baseline_adaptation}).%
\footnote{DS seed-1 LLM-SQL contains the same evaluator-exploiting column-permutation pattern (Case~3, \S\ref{sec:exp:failures}), but only 2 of 5 I2 judges flagged it---below the majority threshold---so it is not counted as a violation in Table~\ref{tab:integrity}. This near-miss illustrates the noise floor of LLM-judged integrity checks at the current vote threshold.}

\paragraph{Reference integrity (I3).}
\textbf{\sys{}} and \textbf{Sakana ASv2} both achieve zero hallucinated references (0/337 and 0/159 respectively).
\textbf{DS} exhibits the highest hallucination rate (42/201, 20.9\%), followed by \textbf{AIR} (21/222, 9.5\%) and \textbf{ARC} (3/196, 1.5\%).
\textbf{ARC}'s low rate reflects its multi-tiered retrieval pipeline (OpenAlex, Semantic Scholar, arXiv, Google Scholar). Its three hallucinated entries are a single fabricated citation (\texttt{sutskever2013importance}, titled ``SGD with Momentum'') from ARC's upstream seminal papers library---a hand-curated YAML file shipped with the framework that assigns an informal title to a real paper (Sutskever et al., ICML 2013, whose actual title is ``On the importance of initialization and momentum in deep learning'').
The entry is injected deterministically into all papers whose topic overlaps with optimization keywords, producing the same fabricated reference in all three EPLB papers.
DS and AIR rely on model memory for reference generation, producing plausible-looking but non-existent bibliography entries---a failure mode illustrated in Case~2 (\S\ref{sec:exp:failures}).
\sys{}'s zero hallucination rate is an architectural property of the PI's citation graph: every reference originates from a Semantic Scholar API call whose result is cached in the evidence chain.
Sakana ASv2's clean record reflects its cached citation retrieval mechanism, which grounds references in API results before paper generation.
(Full list of hallucinated references in Appendix~\ref{appx:discovered_hallucinated_references}.)


\paragraph{Method-code alignment (I4).}
\textbf{\sys{}} achieves 14/15 aligned papers (93\%), compared to \textbf{AIR} (12/15, 80\%), \textbf{Sakana ASv2} (5/15, 33\%), \textbf{DS} (5/15, 33\%), and \textbf{ARC} (3/15, 20\%).
The misalignment patterns differ qualitatively across systems.
\textbf{Sakana ASv2}'s low I4 score is partly attributable to the same design mismatch: the submitted code files contain tuning loops and experiment-tracking code alongside the actual solver, so I4 judges flag this non-solver code as misaligned with the paper's algorithmic claims.
\textbf{AIR}'s failures are typically algorithm mismatches---e.g., the paper describes a more sophisticated procedure than the code implements.
\textbf{ARC} exhibits the worst method-code alignment (3/15, 20\%), a direct consequence of its 23-stage waterfall architecture: code generation (stages~10--13) and paper writing (stages~16--23) run as disconnected phases with no shared intermediate representation.
The paper-writing agent invents algorithm names and describes methods based on experiment metadata, without access to the solver's actual logic, producing algorithm-class mismatches (e.g.\ beam search with Edmonds' arborescence vs.\ greedy edge penalization), undisclosed fallback paths, and method-vs-ablation inversion between paper and submitted code.
\textbf{\sys{}}'s single misaligned paper (\textsc{cloudcast}, 1st seed) is a case where the paper writer fabricated algorithmic claims not present in the code---describing a ``hybrid neuro-symbolic solver'' with ``LLM-guided evolutionary search'' when the submitted code is a deterministic routing heuristic with no LLM calls.
The Claim Verifier's method-code cross-check catches nearly all such misrepresentation before paper finalization.

% \paragraph{Artifact-format caveat.}
% The \coeaudit{} assumes that each system produces a self-contained solver file whose re-execution reproduces the paper's claimed score.
% Sakana ASv2's best-first tree search violates this assumption: its output artifact is the full discovery script---including agent-written evaluator imports for hyperparameter tuning, environment-dependent parameter sweeps, and a framework-injected canonical evaluation harness---rather than an extracted solver function.
% Our code stripping removes the framework harness but not the agent's own tuning infrastructure.
% Sakana's I1 failures reflect a mix of genuine cross-stage cherry-picking (4/7) and artifact-format effects (2/7 environment-dependent tuning, 1/7 metric mismatch); I2 and I4 are confounded by the presence of non-solver code in the audited artifact.
% Cross-system comparison of I2 and I4 should therefore exclude Sakana; I1 and I3 remain valid.
% The remaining systems (ARC, DS, AIR, \sys{}) all produce standalone solver files that satisfy the audit's artifact assumptions.

% We present detailed failure mode case studies in \S\ref{sec:exp:failures} and per-check breakdowns in Appendix~\ref{appx:discovered-score-errors} (I1), \ref{appx:i2_spec_violation} (I2), \ref{appx:discovered_hallucinated_references} (I3), and \ref{app:i4_audit_analysis} (I4).

%\sys{}'s single misaligned paper indicates that the Claim Verifier's method-code cross-check (\S\ref{sec:system:verifier}) catches nearly all misrepresentation before paper finalization, though complete coverage remains a limitation (\S\ref{sec:limitations}).

